% !TeX program = pdflatex
% Self-contained source: pdflatex twice; no fontspec or external bibliography.
\documentclass[11pt,a4paper]{article}
\usepackage[margin=24mm,headheight=14pt]{geometry}
\usepackage[T1]{fontenc}
\usepackage[utf8]{inputenc}
\usepackage{lmodern,microtype}
\usepackage{amsmath,amssymb,mathtools,amsthm}
\usepackage{enumitem,booktabs,longtable,array,xurl,listings,needspace}
\usepackage[hidelinks,unicode]{hyperref}
\hypersetup{pdftitle={Economic Theory: Proof Attempts and Adversarial Checks},pdfauthor={Research notes},pdfsubject={Eleven user-supplied mathematical problems}}
\newcommand{\NE}{\operatorname{NE}}
\DeclareMathOperator*{\argmax}{arg\,max}
\DeclareMathOperator*{\argmin}{arg\,min}
\newtheorem{theorem}{Theorem}[section]
\newtheorem{lemma}[theorem]{Lemma}
\newtheorem{proposition}[theorem]{Proposition}
\newtheorem{corollary}[theorem]{Corollary}
\theoremstyle{definition}
\newtheorem{definition}[theorem]{Definition}
\theoremstyle{remark}
\newtheorem*{remark}{Remark}
\setlength{\parindent}{0pt}
\setlength{\parskip}{5pt plus 1pt minus 1pt}
\setlength{\emergencystretch}{2em}
\setcounter{tocdepth}{1}
\setlist{leftmargin=*,itemsep=2pt,topsep=4pt}
\lstset{basicstyle=\small\ttfamily,columns=fullflexible,keepspaces=true,
 breaklines=true,breakatwhitespace=false,showstringspaces=false,
 frame=single,framerule=0.25pt,aboveskip=8pt,belowskip=8pt,xleftmargin=3pt,xrightmargin=3pt}
\allowdisplaybreaks[1]
\raggedbottom

\begin{document}
{\large\bfseries Research attempt and adversarial verification}\par
9 October 2026\hfill Original-target status: \textbf{UNRESOLVED}\par
\section{C73-27: Large-population selection of potential minimizers}
\label{sec:C73-27}
\noindent\textbf{Input:} \texttt{C73-27.tex}. \quad\textbf{Tools:} web browsing [YES]; Python computation [YES].
\subsection*{1) FORMAL RESTATEMENT}
Fix $d\geq3$, $T>0$, interior $m_0\in\Delta_d$, and smooth potentials $\mathcal F,\mathcal G$ on a neighborhood of the simplex, with $f_i=\partial_i\mathcal F$, $g_i=\partial_i\mathcal G$. A deterministic population control consists of an absolutely continuous simplex path $m$ and measurable rates $a_{ij}\geq0$ ($i\ne j$), with
\[
 \dot m_i=\sum_{j\ne i}(m_ja_{ji}-m_ia_{ij}),\quad m(0)=m_0,
 \quad J(m,a)=\int_0^T\left[\tfrac12\sum_{i,j\ne i}m_i a_{ij}^2+
 \mathcal F(m)\right]dt+\mathcal G(m(T)).
\]
Finite energy means finiteness of the displayed quadratic-rate integral. Let $\mathcal M_{T,m_0}$ be the paths occurring in global minimizers of $J$.

For each integer $N\geq2$, player $\ell$ controls the nonnegative jump rates of its current state, with diagonal rate equal to minus their sum. Its running and terminal costs are
\[
 \tfrac12\sum_{j\ne X_t^\ell}(a_{X_t^\ell j}^\ell(t))^2+
 f_{X_t^\ell}(m_t^{N,-\ell}),\qquad
 g_{X_T^\ell}(m_T^{N,-\ell}),\quad
 m_t^{N,-\ell}=\frac1{N-1}\sum_{r\ne\ell}\delta_{X_t^r}.
\]
Individual transition noises are independent, with no common noise. Initial states are independent with law $m_0$. Equilibria are bounded symmetric Markov-perfect feedbacks: from every time and state vector, no unilateral bounded progressively measurable rate deviation can reduce that player's continuation cost. The bound on a feedback need not be uniform in $N$.

The literal question is whether, for every such datum and every sequence of such equilibria whose empirical laws $\mathcal L(\mu^N)$ converge in Skorokhod $J_1$, with continuous-path limit law $Q$, one has $Q(\mathcal M_{T,m_0})=1$, where $\mu^N=N^{-1}\sum_\ell\delta_{X^\ell}$. The minimal convention $N\geq2$ is necessary because $m^{N,-\ell}$ is undefined at $N=1$; it does not change the limit question. Exact Markov perfection, ex ante approximate Nash, first-order MFG conditions, and global potential optimality are distinct requirements.

\subsection*{2) QUICK LITERATURE/CONTEXT CHECK}
Cardaliaguet and Delarue \cite{cardaliaguet-delarue}, Section 4.3 and the ``Selection of equilibria'' paragraph on printed p.~3691, distinguish the $d\geq3$ large-$N$ selection question from the two-state result and from vanishing-noise selection with corrected costs. The checked source does not justify replacing global minimization by the deterministic equilibrium equations. The present note supplies direct proofs of compactness, an affine-potential solved subclass, and an explicit failure for an ex ante approximate-equilibrium variant. The targeted check was performed on 9 October 2026, not asserted to be an exhaustive survey.

\subsection*{3) ATTACK PLAN}
\textbf{Type and tools.} This is a many-player equilibrium limit with a variational selection requirement. Useful tools are: flux variables (compactness of the cooperative problem); convex duality of quadratic perspectives (lower semicontinuity); Dynkin's formula (single-agent verification); energy bounds from a zero-rate deviation (uniform estimates for exact equilibria); empirical martingales (large-$N$ compactness); Bregman remainders (separate convexity from stationarity); explicit concave terminal potentials (construct nonminimizing equilibria).

\textbf{Proof tracks.} Derive general compactness and an exact global verification identity; solve affine potentials completely. \textbf{Disproof tracks.} Try the zero-control nonminimizing MFG path in a finite game, then test it at every configuration; compare exact and approximate equilibrium quantifiers. The zero profile fails exact Markov perfection, but yields a rigorous counterexample to a weaker ex ante approximate-Nash selection statement.

\subsection*{4) WORK}
\begin{lemma}[Existence and compactness of global minimizing paths]
The deterministic potential problem has a minimizer. Its set $\mathcal M_{T,m_0}$ is compact in the uniform topology on $C([0,T],\Delta_d)$, hence is Borel.
\end{lemma}
\begin{proof}
Set $w_{ij}=m_i a_{ij}\geq0$. The energy becomes
\[
 K(m,w)=\frac12\int_0^T\sum_{i,j\ne i}\frac{w_{ij}^2}{m_i}\,dt,
\]
with $0/0=0$ and $w^2/0=+\infty$ when $w>0$. The path constraint is linear: $\dot m_i=\sum_{j\ne i}(w_{ji}-w_{ij})$. A constant path with zero flux is feasible. Boundedness of the potentials on the simplex implies that a minimizing sequence has uniformly bounded energy. Since $m_i\leq1$, its fluxes are bounded in $L^2$. The forward equation and Cauchy--Schwarz give a uniform $1/2$-H\"older modulus for the paths. Compactness of the simplex and the Arzel\`a--Ascoli theorem yield a uniformly convergent subsequence $m^n\to m$; weak compactness of bounded sets in the Hilbert space $L^2$ yields $w^n\rightharpoonup w$ along a further subsequence. Testing against time-interval indicators passes the integrated forward equation to the limit. Nonnegativity of the fluxes is preserved by weak convergence, since their integrals against every nonnegative $L^2$ test remain nonnegative.

For completeness, lower semicontinuity follows from
\[
 K(m,w)=\sup_{q\in L^\infty}
 \int_0^T\sum_{i,j\ne i}\left[q_{ij}w_{ij}-\tfrac12m_iq_{ij}^2\right]dt.
\]
The pointwise quadratic identity gives ``$\geq$'' for the supremum representation. For the reverse approximation use nonnegative truncated maximizers $q_{ij}=\min\{L,w_{ij}/m_i\}$, assigning $L$ where $m_i=0<w_{ij}$ and zero where both vanish. The integrands increase to the perspective as $L\to\infty$, so monotone convergence gives equality, including the infinite case. For each fixed bounded $q$, its expression is continuous under uniform $m$ and weak $L^2$ flux convergence. Taking their supremum proves the claimed lower semicontinuity.

The potential terms converge by uniform continuity. The limiting pair therefore minimizes. Finite energy forces $w_{ij}=0$ almost everywhere where $m_i=0$; define $a_{ij}=w_{ij}/m_i$ where $m_i>0$, and zero otherwise, to recover an admissible rate representation. The same compactness argument applied to any sequence of minimizers shows that every uniform limit extracted is a minimizing path. Thus $\mathcal M_{T,m_0}$ is compact.
\end{proof}

\begin{theorem}[General compactness of empirical laws of exact equilibria]
For the finite-player games in the input, every sequence of bounded exact Nash feedback profiles (in particular every specified Markov-perfect sequence) has tight empirical path laws in $D([0,T],\Delta_d)$ with $J_1$ topology, and every subsequential limit is concentrated on continuous paths starting at $m_0$.
\end{theorem}
\begin{proof}
Let $f_{\min},f_{\max}$ be the minimum and maximum of all $f_i(m)$ on the compact simplex, and define $g_{\min},g_{\max}$ likewise. A player can deviate to never jumping. Regardless of how the other feedbacks react, that deviation costs at most $Tf_{\max}+g_{\max}$. Exact Nash therefore bounds the player's equilibrium cost by this number. Subtracting the lower bounds on its running and terminal interaction costs gives
\[
 \mathbb E\int_0^T\sum_{j\ne X_t^\ell}
       (a_{X_t^\ell j}^\ell(t))^2dt
 \leq C:=2[T(f_{\max}-f_{\min})+g_{\max}-g_{\min}],
 \tag{E}
\]
uniformly in $\ell,N$. Boundedness of each finite-game rate profile makes all compensators below well defined; (E) supplies the uniform estimates.

Write the empirical semimartingale decomposition as
\[
 \mu^N(t)=\mu^N(0)+\int_0^t B^N(s)ds+M^N(t).
\]
Let $A_N(t)=N^{-1}\sum_\ell\sum_{j\ne X_t^\ell}(a_{X_t^\ell j}^\ell(t))^2$. Each jump contributes opposite signs to two coordinates, so
\[
 \|B^N(t)\|_1\leq\frac2N\sum_{\ell,j\ne X_t^\ell}a_{X_t^\ell j}^\ell(t),
 \qquad \mathbb E\int_0^T\|B^N(t)\|_1^2dt\leq4(d-1)C.
\]
For the second assertion use Cauchy--Schwarz over the $N(d-1)$ current outgoing rates and then (E).

Independent individual jump noises have no simultaneous jumps almost surely. A jump changes the squared Euclidean empirical increment by $2/N^2$. Thus the sum of the coordinate martingale brackets obeys
\[
 \mathbb E\sum_i\langle M_i^N\rangle_T
 \leq\frac2{N^2}\sum_\ell\mathbb E\int_0^T
              \sum_{j\ne X_t^\ell}a_{X_t^\ell j}^\ell(t)dt
 \leq\frac{2\sqrt{T(d-1)C}}N.
\]
Doob's inequality on each coordinate gives
\[
 \mathbb E\sup_{t\leq T}\|M^N(t)\|_2^2
 \leq\frac{8\sqrt{T(d-1)C}}N\longrightarrow0.
\]
Let $Y^N(t)=\mu^N(0)+\int_0^tB^N(s)ds$, a continuous path in $\mathbb R^d$. With arbitrarily high probability its squared-derivative integral is bounded by a deterministic constant, by Markov's inequality and the uniform energy bound. Such paths, with initial points in the compact simplex, form a set with compact closure in the uniform topology: Cauchy--Schwarz bounds increments by a constant times $|t-s|^{1/2}$ and bounds their suprema. Therefore the laws of $Y^N$ are tight on continuous paths. Since $\sup_t\|\mu^N-Y^N\|_2\to0$ in probability, the empirical laws are tight in $J_1$ and every limit is continuous. One can see the latter assertion directly by extracting a weakly convergent subsequence of the tight continuous approximations and using their vanishing uniform discrepancy; uniform distance dominates $J_1$ distance. Finally, the empirical initial law satisfies $\mathbb E\|\mu^N(0)-m_0\|_2^2\leq1/N$. Continuity of time-zero evaluation and the closedness of the simplex imply the stated initial point and path range for every limit.
\end{proof}

\paragraph{What compactness does not establish.}
This theorem supplies the kind of subsequences postulated in the question without assuming a uniform bound on all feedback rates. It does not identify their limiting controls, prove that they solve a deterministic MFG pathwise, or prove global potential optimality. Those are separate steps.

\begin{lemma}[Exact potential verification identity]
Let $(m^*,u)$ be a bounded classical deterministic MFG solution with
\[
 \dot u_i=H_i(u)-f_i(m^*),\quad u_i(T)=g_i(m^*(T)),\quad
 a^*_{ij}=(u_i-u_j)_+,\quad \dot m^*=\operatorname{div}(m^*a^*).
\]
For any admissible finite-energy competitor $(m,a)$ with the same initial law, define
$D_{\mathcal F}(m,m^*)=\mathcal F(m)-\mathcal F(m^*)-f(m^*)\cdot(m-m^*)$, and define $D_{\mathcal G}$ analogously. Then
\begin{align*}
 J(m,a)-J(m^*,a^*)={}&\int_0^T\sum_{i,j\ne i}m_i
 \left[\tfrac12(a_{ij}-(u_i-u_j)_+)^2+(u_j-u_i)_+a_{ij}\right]dt\\
 &+\int_0^TD_{\mathcal F}(m,m^*)dt
 +D_{\mathcal G}(m(T),m^*(T)).
\end{align*}
\end{lemma}
\begin{proof}
Finite energy implies integrable fluxes: $\int_0^Tm_i a_{ij}dt\leq
\sqrt{T\int_0^Tm_i a_{ij}^2dt}$, since $m_i\leq1$. Integration by parts with $u\in C^1$ and the absolutely continuous population paths is therefore justified. The initial difference is zero, and hence
\[
 g(m^*(T))\cdot(m(T)-m^*(T))
 =\int_0^T\left[(H(u)-f(m^*))\cdot(m-m^*)
       +u\cdot(\dot m-\dot m^*)\right]dt.
\]
For a flux $(m_i a_{ij})$, the last linear term is
$u\cdot\dot m=\sum_{i,j\ne i}m_i a_{ij}(u_j-u_i)$.
Expand the potential differences as their linear parts plus $D_{\mathcal F},D_{\mathcal G}$ and substitute this integration-by-parts identity. The terms linear in $f(m^*)$ cancel. Writing $z=u_i-u_j$, the remaining scalar identity is
\[
 \tfrac12a^2-za+\tfrac12z_+^2
 =\tfrac12(a-z_+)^2+(-z)_+a.
\]
For $a=z_+$ the expression is zero. Summing over the competitor fluxes and subtracting the zero expression for $(m^*,a^*)$ proves the displayed identity exactly.
\end{proof}

\begin{proposition}[Convex potentials turn classical equilibrium into global optimality]
If $\mathcal F$ and $\mathcal G$ are convex on the simplex, every bounded classical MFG solution as above is a global potential minimizer, and its population path is the unique minimizing path.
\end{proposition}
\begin{proof}
Convexity implies $D_{\mathcal F},D_{\mathcal G}\geq0$ by the first-order convexity inequality on each line segment in the simplex. Every other term in the identity is nonnegative because rates and masses are nonnegative. Thus $J(m,a)\geq J(m^*,a^*)$ for every admissible competitor. Equality forces $m_i(a_{ij}-a^*_{ij})^2=0$ almost everywhere, and hence $m_i a_{ij}=m_i a^*_{ij}$ almost everywhere. The competing population then solves the same bounded linear forward ODE with rates $a^*(t)$ and the same initial law. Uniqueness follows from Gronwall's inequality applied to the difference, so $m=m^*$. No strict convexity of the potentials was needed. This proposition is conditional on identifying a classical MFG limit; it is not such an identification theorem for all finite-game limits.
\end{proof}

\begin{theorem}[Full selection in the affine-potential subclass]
Suppose $\mathcal F(m)=c\cdot m+C_0$ and $\mathcal G(m)=\gamma\cdot m+C_1$. Every sequence of equilibria from the input has empirical paths converging uniformly in probability to the unique global minimizing path $m^*$. In particular every subsequential law limit is $\delta_{m^*}$ and satisfies the requested selection property.
\end{theorem}
\begin{proof}
Here the individual interaction costs are $f_i=c_i$ and $g_i=\gamma_i$, independent of all other players' states. Solve the backward ODE
\[
 \dot u_i=H_i(u)-c_i,\qquad u(T)=\gamma,
 \qquad a^*_{ij}(t)=(u_i(t)-u_j(t))_+.
\]
Local existence and uniqueness hold because $H$ is locally Lipschitz. In reversed time $v(\tau)=u(T-\tau)$ the equation is $v_i'=c_i-H_i(v)$. A maximal component has derivative at most $\|c\|_\infty$; a minimal component has $H_i=0$ and derivative at least $-\|c\|_\infty$. These assertions also hold as one-sided derivatives at ties, by choosing an active extremal index. Equivalently, comparison with the scalar barriers $\pm(\|\gamma\|_\infty+\tau\|c\|_\infty)$ gives
$\|u(t)\|_\infty\leq M:=\|\gamma\|_\infty+T\|c\|_\infty$.
This prevents finite-time escape and extends the ODE through the whole horizon. Its feedback rates are continuous and at most $2M$.

For a single controlled state process $X$, possibly with rates depending on all players' observed histories, apply Dynkin's formula to $u_i(t)$ at $i=X_t$. For every bounded admissible control its expected cost equals
\[
 u_{X_0}(0)+\mathbb E\int_0^T\sum_{j\ne X_t}
 \left[\tfrac12(a_{X_tj}-a^*_{X_tj})^2+
              (u_j-u_{X_t})_+a_{X_tj}\right]dt.
\]
This is the same scalar square-completion as in the preceding identity; the ODE and terminal condition cancel the other terms. Thus $a^*$ attains the minimum against any behavior of the other players. A Nash player's actual cost must attain this minimum, since it may deviate to $a^*$. The nonnegative square remainder then vanishes almost everywhere with respect to time and its realized law, forcing its actual outgoing intensities to equal $a^*$ on that law.

Consequently the joint state process has, up to null time--path sets, the deterministic generator equal to the sum of the independent single-particle generators with rates $a^*(t)$. A finite-state time-inhomogeneous generator with bounded continuous coefficients determines its transition matrices uniquely by the linear Kolmogorov equations. Their product transition law is a solution, hence is the unique joint law. The initial distribution is a product as well. This justifies independence of the players' realized states in this subclass, even if an equilibrium feedback was written using irrelevant other-player coordinates.

Let $m^*$ solve the corresponding linear forward equation from $m_0$. In the empirical decomposition the drift is the linear map $b(\mu^N,a^*)$. Its Lipschitz constant is bounded by some $L$ depending only on $d,M$. The coordinate martingale brackets are at most a constant times $T/N$, because each empirical jump has size $1/N$ and the total jump intensity is at most $2NM(d-1)$. Doob's inequality, the initial bound $\mathbb E\|\mu^N(0)-m_0\|_2^2\leq1/N$, and the pathwise integral inequality give
\[
 \sup_{t\leq T}\|\mu^N(t)-m^*(t)\|_2
 \leq e^{LT}\left(\|\mu^N(0)-m_0\|_2+
                 \sup_{t\leq T}\|M^N(t)\|_2\right),
\]
\[
 \mathbb E\sup_{t\leq T}\|\mu^N(t)-m^*(t)\|_2^2\leq\frac{C_{d,T,c,\gamma}}N.
\]
This proves uniform convergence in probability. Finally affine potentials have zero Bregman remainders, and $(m^*,u)$ satisfies the classical MFG equations. The verification identity proves global optimality and uniqueness of its minimizing population path. The argument applies to every specified equilibrium, not only a selected favorable one.
\end{proof}

\begin{proposition}[An explicit nonminimizing MFG and a failed finite-game candidate]
For $d=3$, $T=1$, $m_0=(1/3,1/3,1/3)$, $\mathcal F=0$ and
$\mathcal G(m)=-3(m_1-1/3)^2$, the constant zero-control MFG path is not a potential minimizer. The profile prescribing zero rates in every finite-player configuration is not Markov perfect for any $N\geq2$.
\end{proposition}
\begin{proof}
The constant law and values $u=0$ satisfy the deterministic backward--forward equations because $g(m_0)=0$. Its objective is zero. The control $a_{21}=a_{31}=\log2$, all other rates zero, gives
$m(t)=(1-\frac23e^{-t\log2},\frac13e^{-t\log2},\frac13e^{-t\log2})$ and
$J=(\log2)/6-1/3<0$. This verifies nonminimality with an explicit finite-energy competitor.

To check the finite game, take a configuration where all players other than $\ell$ are in state one and player $\ell$ is in state two. Under zero feedbacks its continuation cost is zero. Its opponents remain in state one, so its own terminal costs are $g_1=-4,g_2=g_3=0$. Deviate by jumping from state two to one at constant rate $c\in(0,8)$ and making no further jumps. Over a remaining horizon $r>0$ this has expected cost
\[
 \left(\tfrac c2-4\right)(1-e^{-cr})<0.
\]
The energy term is $\frac12c^2\int_0^r e^{-ct}dt$, and the terminal term is $-4(1-e^{-cr})$, yielding the formula. The deviation is bounded and admissible. Thus zero feedback fails the required continuation Nash inequality. In fact this configuration has positive probability under the prescribed independent interior initial law, however small that probability becomes as $N$ grows.
\end{proof}

\begin{proposition}[A complete counterexample for the ex ante approximate-Nash variant]
In the same three-state example, the zero-feedback profiles are ex ante $\eta_N$-Nash equilibria with $\eta_N=2\sqrt{2/(N-1)}\to0$. Their empirical laws converge to the nonminimizing constant MFG path. Thus replacing exact Markov perfection by ex ante vanishing-error Nash would make the selection assertion false.
\end{proposition}
\begin{proof}
Fix a player. When all other players have zero rates, their initial empirical fraction $Y$ in state one is constant in time, with mean $1/3$ and variance $2/[9(N-1)]$. Its own terminal cost in state one is $g_1(Y)=-6(Y-1/3)$ and in the other states it is zero. For any bounded progressively measurable deviation, kinetic cost is nonnegative and terminal cost is at least $-|g_1(Y)|$. Therefore its deviating expected cost is at least
\[
 -\mathbb E|g_1(Y)|\geq-6\sqrt{\operatorname{Var}(Y)}
 =-2\sqrt{2/(N-1)}.
\]
Under zero rates, its initial state is independent of $Y$, so its expected cost is $(1/3)\mathbb E g_1(Y)=0$. This bounds every ex ante unilateral improvement by $\eta_N$. Every empirical path is its initial empirical law held constant, which converges in probability to $m_0$. Nonminimality was proved above. This approximate result does not hold uniformly over continuation configurations, as the preceding deviation with $g_1=-4$ demonstrates.
\end{proof}

\paragraph{Computation.}
The verification script checks the explicit competitor's terminal law $(2/3,1/6,1/6)$ and objective approximately $-0.2178088032$. The square-completion identity is exact algebra, and the finite-game deviation inequality follows from $0<c<8$, not a numerical optimization. Minimal pseudocode for testing the proposed zero finite-game feedback is:
\begin{lstlisting}
for N in [2, 3, 4]:
    # Other N-1 players all occupy state 1; deviator starts in state 2.
    for c in [1, 2, 4]:
        deviation_cost = (c/2 - 4)*(1-exp(-c))
        assert deviation_cost < 0
\end{lstlisting}
The supplied script evaluates the same analytic construction; neither these cases nor a simulation is used to assert an exact-Nash counterexample.

\subsection*{5) VERIFICATION}
The flux formulation proves that the event $\mathcal M_{T,m_0}$ is a Borel event, avoiding a hidden measurability issue in $Q(\mathcal M)=1$. The compactness theorem uses an actual allowed zero-rate deviation and independent jump compensators, not a uniform rate bound absent from the input. The verification identity includes the extra term $(u_j-u_i)_+a_{ij}$; omitting it when $u_i<u_j$ would make the square completion false. Finite energy suffices for all integrations by parts. The affine theorem covers all equilibria up to irrelevant null-set prescriptions. The nonconvex candidate was explicitly tested at rare off-path configurations and fails exact Markov perfection. Its valid approximate-Nash counterexample is labeled with that weaker, ex ante quantifier only.

\Needspace{8\baselineskip}
\subsection*{6) FINAL}
\textbf{LABEL: UNRESOLVED} for the exact Markov-perfect selection question with arbitrary smooth nonconvex potentials.

\textbf{Strongest checked partial results.} General compactness of exact-equilibrium empirical laws with continuous limits; compactness of the global minimizing-path set; an exact potential verification identity; full selection for every equilibrium in the affine-potential subclass; and a complete counterexample to the weaker ex ante approximate-Nash formulation.

\textbf{Exact first gap.} Exact finite-player Nash inequalities have not been converted into a global potential-minimization inequality for every limiting population path when the potential Bregman remainders can be negative.

\textbf{Top three next moves.} (1) Quantify the discrepancy between the finite-game Nash system and a discrete global potential optimization problem, including all $m^{N,-\ell}$ shifts. (2) Identify limit controls using the proved energy compactness while retaining the continuation optimality information, not merely ex ante approximate optimality. (3) Test the explicit concave terminal potential using actual finite-$N$ Markov-perfect feedbacks rather than the rejected zero-rate candidate.

\textbf{Minimal counterexample perspective.} A candidate must retain a nonminimizing path while remaining an exact best response at every finite configuration, including exponentially rare ones. The displayed zero-rate profiles show why merely approximating a deterministic nonminimizing equilibrium is insufficient.

\paragraph{Reproducibility and scope.} This is one of eleven companion reports. The bundle includes \texttt{verify.py} and its executed results. Finite experiments supplement the proofs; they are not proofs of asymptotic claims. Literature coverage is targeted, and no novelty or formal proof-assistant certification is claimed.

\begin{thebibliography}{99}
\small
\bibitem{cardaliaguet-delarue}
Pierre Cardaliaguet and Fran\c{c}ois Delarue.
\emph{Selected topics in mean field games}. Proceedings of the International Congress of Mathematicians 2022, volume 5.
DOI: 10.4171/ICM2022/17. Sections 4.3--4.4; especially printed pp.~3691--3693.
\url{https://ems.press/content/book-chapter-files/33265}.

\end{thebibliography}

\end{document}
